\section{Formal and computational verification}\label{app:verification}
\begingroup
\def\UrlBreaks{\do\.}\def\UrlBigBreaks{\do\/}
\subsection{Formal results}
The negative conclusion of Theorem~\ref{thm:negative} is formally proved
in Lean~4 with Mathlib, unconditionally. The formal project defines the literal quartic as
\path|ToeplitzSOS.toeplitzBW| and uses the predicate
\path|ToeplitzSOS.IsSumSqHomQuad| for finite sums of real homogeneous
quadratic squares. The theorem
\path|ToeplitzSOS.Negative.sharpNegativeResolution| proves
\path|SharpNegativeResolution|, namely non-SOS for every
$N\ge\ObstructionOrder$. The declarations
\path|ToeplitzSOS.Negative.sharp_negative_explicit| and
\path|ToeplitzSOS.Negative.not_isSumSqHomQuad_of_large_order| give the
expanded quadratic quantifiers and the direct predicate form,
respectively. The declaration
\path|ToeplitzSOS.Negative.not_MI15| refutes the registered catalogue
statement. The classical nonnegativity assertion is proved in
Section~\ref{sec:context}. The formal positive range is $2\le N\le20$.


The proof composes
\path|ToeplitzSOS.Negative.cornerGram_of_sos|, which extracts the complete
corner data from an arbitrary SOS of the literal polynomial, with
\path|ToeplitzSOS.Negative.no_cornerGram|, which excludes those data at
the selected depth. The analytic inputs are supplied by the proved
\path|ToeplitzSOS.Negative.Analytic.analyticInputs|. They include
\path|ToeplitzSOS.Negative.Analytic.baselineTail|, the mixed correction bound, and
the two-rectangle comparison; their coefficient, complete-mean, and
majorant hypotheses are proved for the corner data extracted from an
arbitrary sum of squares.
The integer witness is checked by reduction in the Lean kernel in
\path|ToeplitzSOS.Negative.WitnessCertificate|, and
\path|ToeplitzSOS.Negative.Witness.tangent_pairing_lt| proves its negative
kernel pairing.

The formal final budget uses the exact witness mass bound
$\|c\|_1^2\le3\cdot2^{58}$. At the selected scale the realigned entries
are at most $2^{-18}$, so their contribution is at most
$3\cdot2^{58}\cdot2^{-18}=3\cdot2^{40}<2^{42}$; the finite-to-tangent
contribution is at most $3/4$. Thus the total error is less than
$2^{42}$ without the convexity step used for the coarser mass bound in
Corollary~\ref{cor:threshold}.

The formal proofs use only the audited standard axioms
\path|propext|, \path|Classical.choice|, and \path|Quot.sound|.

\subsection{The result map}
The table lists the numbered results; declaration names are relative to
the namespace \path|ToeplitzSOS|. The last column names the exact
computation behind a result, identified in the accompanying
verification map; ``Written proof'' means that the argument in this
paper uses no computation. The Lean declarations give formal proofs
within the scopes stated in the table. Some intermediate results are
organized differently in the formal proof: it uses polarization in
place of the full relation-space classification and closes the finite
contradiction at the selected scale with the budget just described.
The exact positive chain proves orders through $50$, and its Lean scope
is as shown; rational Gram certificates suffice for the real problem
(Section~\ref{sec:positive}).

\begingroup
\footnotesize
\setlength{\tabcolsep}{3pt}
\def\UrlBreaks{\do\.}\def\UrlBigBreaks{\do\/}
\begin{longtable}{@{}>{\raggedright\arraybackslash}p{.33\textwidth}>{\raggedright\arraybackslash}p{.43\textwidth}>{\raggedright\arraybackslash}p{.21\textwidth}@{}}
\toprule Result & Lean declaration & Exact check\\
\midrule\endhead
Theorem~\ref{thm:negative} Negative answer & {\scriptsize\path|Negative.sharpNegativeResolution|}\newline {\scriptsize\path|Negative.sharp_negative_explicit|}\newline Unconditional non-SOS assertion; nonnegativity is classical. & \path|negative.py|\newline \path|witness.py|\\[5pt]
Theorem~\ref{thm:finite-scale} Uniform finite-scale obstruction & {\scriptsize\path|Negative.finite_separation|}\newline {\scriptsize\path|Negative.no_feasible_evaluation|}\newline Formal separation at the selected scale; the uniform budget is the written version. & \path|negative.py|\newline \path|witness.py|\\[5pt]
Corollary~\ref{cor:threshold} Explicit threshold & {\scriptsize\path|Negative.cornerGram_impossible|}\newline {\scriptsize\path|Negative.not_isSumSqHomQuad_of_large_order|}\newline The selected corner and every larger ambient order. & \path|negative.py|\newline \path|witness.py|\\[5pt]
Lemma~\ref{lem:frame} The complete quadratic frame & {\scriptsize\path|Negative.exists_alternating_sos|}\newline {\scriptsize\path|Negative.alternating_eq_wedge_sum|}\newline SOS-to-wedge reduction; the full relation-space classification is replaced by polarization. & Written proof\\[5pt]
Lemma~\ref{lem:corner} Stabilization and the complete normal form & {\scriptsize\path|Negative.outerCorner_eq_baselinePolynomial|}\newline {\scriptsize\path|Negative.cornerGram_of_sos|}\newline Complete literal corner extraction from arbitrary SOS data. & Written proof\\[5pt]
Lemma~\ref{lem:means} Forced complete means & {\scriptsize\path|Negative.corner_sos_low_means|}\newline {\scriptsize\path|Negative.baselineMixedEntry_low_column|}\newline Complete low means; baseline positivity is given by the written Laplacian proof. & Written proof\\[5pt]
Lemma~\ref{lem:global} Global bound and the repeated-node identity & {\scriptsize\path|Negative.CornerGram.repeated_diagonal|}\newline {\scriptsize\path|Negative.referenceQ_conjugate|}\newline {\scriptsize\path|Negative.CornerGram.plusVector_global|}\newline Kernel identities and a sufficient global majorant; the formal radial estimate uses a looser constant. & Written proof\\[5pt]
Lemma~\ref{lem:tails} Uniform finite tails & {\scriptsize\path|Negative.Analytic.baselineTail|}\newline {\scriptsize\path|Negative.CornerGram.repeated_tail_bound|}\newline Both baseline truncations and the combined repeated-node error; the corner extraction supplies the coefficient and mean hypotheses. & Written proof\\[5pt]
Lemma~\ref{lem:continuation} Uniform improvement of the pure kernel & {\scriptsize\path|Negative.Analytic.twoRectangles|}\newline {\scriptsize\path|Negative.CornerGram.realigned_bound|}\newline General comparison; Gram majorants and the realigned witness entries are discharged at the selected scale. & Written proof\\[5pt]
Lemma~\ref{lem:witness} Exact negative pairing & {\scriptsize\path|Negative.Witness.tangent_pairing_lt|}\newline {\scriptsize\path|Negative.Witness.fullUpper_lt|}\newline Exact integer certificate checked by Lean kernel reduction. & \path|negative.py|\newline \path|witness.py|\\[5pt]
Proposition~\ref{prop:positive} Exact positive certificates & {\scriptsize\path|toeplitzBW_isSumSq_of_le_20|}\newline Orders 2 through 20 only. & \path|positive.py|\\[5pt]
Theorem~\ref{thm:exchange} Exact exchange criterion & Not formalized & Written proof\\[5pt]
Theorem~\ref{thm:criticality} Criticality of the pure budget & Not formalized & Written proof\\[5pt]
Corollary~\ref{cor:boundary} Failure of the all-order increment condition & Not formalized & Written proof\\[5pt]
Remark~\ref{rem:persistence} An expectation at moderate orders & Not formalized & None (unproved remark)\\[5pt]
\bottomrule
\end{longtable}
\endgroup


\endgroup
